\section{Hierarchical sampling and exchangeability}\label{sec::hier_sampling}

% \textcolor{blue}{The hierarchical sampling procedure: samples of samples. Put examples where it has already been used. And it actually corresponds to iid samples of the first terms of an exchangeable sequence. Definition of the hierarchical empirical measure (HEM).}


% polished start
We assume that the space $\X$ on which probability measures are defined is a Polish space (i.e., a complete, separable metric space) with a bounded diameter, meaning $\operatorname{diam}(\X) = \sup_{x,y \in \X} d_{\X}(x,y) < \infty$. Let $\probx$ denote the space of probability measures on $\X$. We will work with Random Probability Measures (RPMs), which are random variables taking values in $\probx$. The laws of RPMs are objects in $\probprobx$, which denotes the space of probability measures on $\probx$. Unless stated otherwise, the $\sigma$-algebra on $\probprobx$ is the Borel $\sigma$-algebra generated by the topology of weak convergence, where the continuous functions on $\probx$ are defined through the weak convergence on $\probx$. 

In practice, it is rare to observe the probability measures themselves directly, especially when dealing with continuous probability distributions. Instead, one obtains observations from each latent probability measure. Therefore, we assume that given a law of a random probability measure $Q \in \probprobx$, instead of observing $\widetilde{P}_1, \dots, \widetilde{P}_n$ from $Q$, we observe samples from each $\widetilde{P}_i$. Such a collection of samples is referred to as a hierarchical sample:

\begin{definition}[Hierarchical sample]
    \label{def:hier_sample}
    Given $Q \in \probprobx$, the hierarchical sample refers to the collection of random variables $\{X_{i,j} : i=1,\dots,n,\ j=1,\dots,m\}$, sampled as 
    \[
    \begin{pmatrix}
    X_{1,1} & X_{1,2} & \cdots & X_{1,m} \mid \widetilde{P}_1 \overset{\mathrm{i.i.d.}}{\sim} \widetilde{P}_1 \\
    X_{2,1} & X_{2,2} & \cdots & X_{2,m} \mid \widetilde{P}_2 \overset{\mathrm{i.i.d.}}{\sim} \widetilde{P}_2 \\
    \vdots & \vdots & \ddots & \vdots \\
    X_{n,1} & X_{n,2} & \cdots & X_{n,m} \mid \widetilde{P}_n \overset{\mathrm{i.i.d.}}{\sim} \widetilde{P}_n
    \end{pmatrix},
    \] 
    where $\widetilde{P}_1, \dots, \widetilde{P}_n \overset{\mathrm{i.i.d.}}{\sim} Q$.
\end{definition}

Throughout this paper, $n$ and $m$ will refer to the number of rows and columns of the hierarchical sample, respectively. We denote the hierarchical sample by $\mathbf{X}_{n,m} = (X_{i,j})$, where boldface letters refer to matrices and non-bold letters refer to their individual elements. With a slight abuse of notation, we will also sometimes denote $\mathbf{X}_{n,m} \sim Q$. Hierarchical samples are what we observe to decide whether two laws of random probability measures are the same. Before proceeding, let us discuss the limitations of the hierarchical sample.

Let $Q^1, Q^2 \in \probprobx$. Suppose that we observe two hierarchical samples $\mathbf{X}_{n,m} \sim Q^1$ and $\mathbf{Y}_{n,m} \sim Q^2$. Note that each of the $n$ sequences of length $m$ from the first hierarchical sample consists of i.i.d. random variables from the following law on $\X^m$: for $i = 1, \dots, n$,
$$
P^{1,(m)}(A_1 \times \cdots \times A_m) := \mathbb{P}(X_{i,1} \in A_1, \dots, X_{i,m} \in A_m) = \int_{\probx} \prod_{j=1}^m P(A_j) Q^1(dP).
$$
Similarly, for the second hierarchical sample, we have
$$
P^{2,(m)}(A_1 \times \cdots \times A_m) := \mathbb{P}(Y_{i,1} \in A_1, \dots, Y_{i,m} \in A_m) = \int_{\probx} \prod_{j=1}^m P(A_j) Q^2(dP).
$$
Such a characterization implies that the sequences $X_{i,1}, \dots, X_{i,m}$ and $Y_{i,1}, \dots, Y_{i,m}$ are finitely exchangeable. This means that each of the hierarchical samples, $\mathbf{X}_{n,m}$ and $\mathbf{Y}_{n,m}$, contains $n$ finitely exchangeable sequences that are independent of one another. It is a well-known fact in Bayesian nonparametrics that a finite exchangeable sequence is not always in one-to-one correspondence with the law of a random probability measure. An implication of this is that it might happen that the laws of the RPMs are different, yet the hierarchical samples we observe are statistically indistinguishable; that is, all the finite exchangeable sequences we observe have the same probability distribution.

Indeed, consider two distinct Dirichlet processes 
\[
    Q^1 = DP(\alpha_1, P_0), \quad Q^2 = DP(\alpha_2, P_0),
\]
with $\alpha_1 \neq \alpha_2$. Then $Q^1 \neq Q^2$. Suppose that the size of the hierarchical samples from $Q^1$ and $Q^2$ is $n \times 1$. Then, the laws characterizing the two hierarchical samples are the same; i.e., for all $A \in \mathcal{B}(\X)$,
\begin{align*}
&P^{1,(1)}(A) = \mathbb{P}(X_{i,1} \in A) = \int_{\probx} P(A) \, Q^1(dP), \\
&P^{2,(1)}(A) = \mathbb{P}(Y_{i,1} \in A) = \int_{\probx} P(A) \, Q^2(dP) =  P_0(A). 
\end{align*}
In this case, no testing procedure can achieve a low Type II error, which is the probability of failing to reject the null hypothesis when it is false. To overcome this limitation, we must ensure that 
\[
Q^1 = Q^2 \quad \text{if and only if} \quad P^{1,(m)} = P^{2,(m)}.
\]
This indeed holds true for any $Q^1, Q^2$ for a sufficiently large $m$ (see Lemma \ref{lemma::finite_exchseq_equivalent_infinite_exchseq}).

As noted above, hierarchical samples $\mathbf{X}_{n,m} \sim Q^1$ and $\mathbf{Y}_{n,m} \sim Q^2$ are used to decide whether $Q^1 = Q^2$. We aim to develop a testing scheme inspired by the Kolmogorov-Smirnov (KS) test. Recall that in the KS test, one rejects the null hypothesis if the distance between the empirical measures of the samples exceeds a given threshold. In Section \ref{sec::distances}, we will equip the space $\probprobx$ with a distance function; for the moment, assume that $d$ defines a distance on that space. Then, we could naturally extend the idea of the KS test to our setting if we had access to the samples $\widetilde{P}^1_1, \dots, \widetilde{P}^1_n \overset{\mathrm{i.i.d.}}{\sim} Q^1$ and $\widetilde{P}^2_1, \dots, \widetilde{P}^2_n \overset{\mathrm{i.i.d.}}{\sim} Q^2$, by rejecting $H_0$ if 
\[
d\left( \frac{1}{n}\sum_{i=1}^n \delta_{\widetilde{P}^1_i}, \frac{1}{n}\sum_{i=1}^n \delta_{\widetilde{P}^2_i}\right)
\]
exceeds some threshold. However, recall that we do not observe $\widetilde{P}^1_i$ and $\widetilde{P}^2_i$. Nevertheless, we can estimate them from the observations they generated; that is, we estimate each of them using the corresponding sequences $X_{i,1}, \dots, X_{i,m}$ and $Y_{i,1}, \dots, Y_{i,m}$ as
\[
\widehat{P}^1_{i,m} = \frac{1}{m}\sum_{j=1}^m \delta_{X_{i,j}}, \quad \widehat{P}^2_{i,m} = \frac{1}{m}\sum_{j=1}^m \delta_{Y_{i,j}}, \quad i = 1, \dots, n.
\]
Thus, the idea of the test is to reject $H_0$ if 
\[
d\left( \frac{1}{n}\sum_{i=1}^n \delta_{\widehat{P}^1_{i,m}}, \frac{1}{n}\sum_{i=1}^n \delta_{\widehat{P}^2_{i,m}}\right)
\]
is higher than some threshold.

Note that the law of each $\widehat{P}^k_{i,m}$ is different from $Q^k$ for $k = 1,2$. Since these random probability measures are the main objects of study in this paper, we introduce a specific notation for their laws. 

In particular, given $Q \in \probprobx$, we denote by $Q^{(m)}$ the law of the empirical measure constructed from an exchangeable sequence; i.e.,
\[
Q^{(m)} = \law \left( \frac{1}{m}\sum_{j = 1}^m \delta_{Z_j}\right),
\]
where $Z_1, \dots, Z_m$ is an exchangeable sequence generated as
\begin{align*}
P &\sim Q, \\
Z_1, \dots, Z_m \mid P &\overset{\mathrm{i.i.d.}}{\sim} P.
\end{align*}

To summarize, given $Q^1$ and $Q^2$, since we estimate the random probability measures $\widehat{P}^k_{i,m}$ from the hierarchical samples $\mathbf{X}_{n,m}$ and $\mathbf{Y}_{n,m}$, we effectively observe 
\[
\widehat{P}^1_{1,m}, \dots, \widehat{P}^1_{n,m} \overset{\mathrm{i.i.d.}}{\sim} Q^{1,(m)}, \quad \widehat{P}^2_{1,m}, \dots, \widehat{P}^2_{n,m} \overset{\mathrm{i.i.d.}}{\sim} Q^{2,(m)}.
\]
We aim to determine whether $Q^1$ and $Q^2$ are identical by comparing the empirical measures concentrated on these two samples. These empirical measures are formally defined as follows:

\begin{definition}[Hierarchical estimator]
    \label{def:hier_estimator}
    Given $Q \in \probprobx$, and a hierarchical sample $\mathbf{X}_{n,m}$, the hierarchical estimator of $Q$ is defined as 
    \[
    Q_{n,m} = \frac{1}{n}\sum_{i = 1}^n \delta_{\widehat{P}_{i,m}},
    \]
    where $\widehat{P}_{i,m} = \frac{1}{m}\sum_{j = 1}^m \delta_{X_{i,j}}$ for all $i = 1, \dots, n$.
\end{definition}

The testing scheme we develop will reject the null hypothesis whenever the hierarchical estimators $Q^1_{n,m}$ and $Q^2_{n,m}$ are far apart. To clarify what is meant by "far apart," we require a distance function on $\probprobx$.

%polished end



% We assume that the space $\obsspace$ on which probability measures are defined is a Polish space (i.e., a complete, separable metric space) with bounded diameter, meaning $\operatorname{diam}(\X) = \sup_{x,y \in \X} d_{\X}(x,y) < \infty$. Let $\lawsspace$ denote the space of probability measures on $\X$. We will be working with the Random Probability Measures (RPMs), which are random variables taking values in $\probx.$  Laws of RPMs are the objects in $\probprobx$, which denotes the space of probability measures on $\probx$. Unless stated otherwise, the sigma-algebra on $\probprobx$ is the Borel sigma-algebra generated by the topology of weak convergence, where the continuous functions on $\probx$ are defined through the weak convergence on $\probx$. 

% In practice it is rare to directly observe the probability measures themselves, especially if we are dealing with continuous probability distributions. Instead, one obtains the observations from each latent probability measure. Therefore, we assume that given a law of random probability measure $Q\in\probprobx$, instead of observing $\widetilde{P}_1,\ldots, \widetilde{P}_n$ from $Q$, we observe samples from each $\widetilde{P}_i.$ Such a collection of samples is referred to as a hierarchical sample:
% \begin{definition}[Hierarchical sample]
%     \label{def:hier_sample}
%     Given $Q\in\probprobx,$ the hierarchical sample refers to the collection of random variables $\left\{X_{i,j} : i=1,\dots,n,\ j=1,\dots,m\right\}$, sampled as 
%     \[
%     \begin{pmatrix}
%     X_{1,1} & X_{1,2} & \cdots & X_{1,m} | \widetilde{P}_1 \overset{\mathrm{i.i.d.}}{\sim} \widetilde{P}_1\\
%     X_{2,1} & X_{2,2} & \cdots & X_{2,m} | \widetilde{P}_2 \overset{\mathrm{i.i.d.}}{\sim} \widetilde{P}_2 \\
%     \vdots & \vdots & \ddots & \vdots \\
%     X_{n,1} & X_{n,2} & \cdots & X_{n,m} | \widetilde{P}_n \overset{\mathrm{i.i.d.}}{\sim} \widetilde{P}_n
%     \end{pmatrix},
%     \] where $\widetilde{P}_1,\ldots,\widetilde{P}_n\overset{\mathrm{i.i.d.}}{\sim} Q.$
% \end{definition}
% Throughout this paper, $n$ and $m$ will refer to the number of rows and columns of the hierarchical sample, respectively. We denote the hierarchical sample by $\mathbf{X}_{n,m} = (X_{i,j})$, where boldface letters refer to matrices and non-bold letters refer to their individual elements. With the abuse of notation, we will also sometimes denote $\mathbf{X}_{n,m} \sim Q.$ Hierarchical samples are what we observe to decide whether two laws of random probability measures are the same. Before proceeding, let us discuss the limitations of the hierarchical sample. 

% Let $Q^1, Q^2 \in \probprobx.$ Suppose that we observe two hierarchical samples $\mathbf{X}_{n,m} \sim Q^1, \mathbf{Y}_{n,m}\sim Q^2.$ Note that each of the $n$ sequences of length $m$ from the first hierarchical sample are i.i.d random variables from the following law on $\X^m$: for $i = 1,\cdots n,$
% $$
% P^{1,(m)}\left(A_1\times \cdots \times A_m\right) := \mathbb{P}(X_{i,1}\in A_1,\ldots,X_{i,m}\in A_m) = \int_{\probx} \prod_{j = 1}^m P(A_j)Q^1(dP).
% $$
% Similarly, for the second hierarchical sample we have
% $$
% P^{2,(m)}\left(A_1\times \cdots \times A_m\right) :=\mathbb{P}(Y_{i,1}\in A_1,\ldots,Y_{i,m}\in A_m) = \int_{\probx} \prod_{j = 1}^m P(A_j)Q^2(dP).
% $$
% Such a characterization implies that the sequences $X_{i,1},\ldots, X_{i,m}$ and $Y_{i,1},\ldots, Y_{i,m}$ are finitely exchangeable. This means that each of the hierarchical samples, $\mathbf{X}_{n,m}$ and $\mathbf{Y}_{n,m},$ contains $n$ finitely exchangeable sequences which are independent across each other. It is a well-known fact in Bayesian nonparametrics that a finite exchangeable sequence is not always in one-to-one correspondence with the law of a random probability measure. Implication of this is that it might happen that laws of RPMs are different, but the hierarchical samples we observe are statistically indistinguishable, meaning that, all of the finite exchangeable sequences we observe have the same probability distribution.

% Indeed, consider two distinct Dirichlet processes 
% \[
%     Q^1 = DP(\alpha_1, P_0), \quad Q^2 = DP(\alpha_2, P_0),
% \]
% with $\alpha_1 \neq \alpha_2$. Then $Q^1 \neq Q^2$. Suppose that size of the hierarchical samples from $Q^1$ and $Q^2$ is $n \times 1.$ Then, the laws characterizing two hierarhical samples are the same, i.e. for all  $A \in  \mathcal{B}(\X),$
% \[
%     P^{1,(1)}\left(A\right) = \mathbb{P} \left(X_{i,1} \in A \right) = \int_{\probx} P(A) \, Q^1(dP) 
%     = P_0(A) = \mathbb{P} \left(Y_{i,1} \in A \right)=P^{2,(1)}\left(A\right).
% \]
% By de Finetti's theorem, $Q$ is in one-to-one correspondence to the law of an infinite sequence of exchangeable random variables, denoted by $P^\infty,$ which is a probability measure on $\X^\infty.$ Its law is characterized by: 
% $$
% P^\infty(A_1\times...\times A_s\times \X\times ...\times \X)=\int_{\probx} \prod_{j = 1}^s P(A_j)Q(dP) \qquad s\geq 1, A_1,\ldots,A_s\in\mathcal{B}(\X).
% $$
% Therefore, the hierarchical sampling scheme can be viewed as the sampling finite sequences from $P^\infty.$ In particular, we observe the i.i.d. samples from the pushforward measure $P^\infty\circ \pi_m^{-1},$ defined by:
% $$
% P^\infty\circ \pi_m^{-1}(A) = P^\infty(\pi_m^{-1}(A)) \qquad A \in \mathcal{B}(\X^m),
% $$ where $\pi_m : \X^\infty \to \X^m$ denotes the projection onto the first $m$ coordinates.
% \begin{remark}
%     There exist $Q^1, Q^2 \in \probprobx$ with $Q^1 \neq Q^2$, that is, $P_1^\infty \neq P_2^\infty$, while still satisfying $P_1^\infty \circ \pi_m^{-1} = P_2^\infty \circ \pi_m^{-1}$. 

%     Indeed, consider two distinct Dirichlet processes 
%     \[
%         Q^1 = DP(\alpha_1, P_0), \quad Q^2 = DP(\alpha_2, P_0),
%     \]
%     with $\alpha_1 \neq \alpha_2$. Then $Q^1 \neq Q^2$, and consequently their induced laws on infinite exchangeable sequences, $P_1^\infty$ and $P_2^\infty$, are also distinct.  

%     Now let $m = 1$ and take any $A \in \mathcal{B}(\X)$. We compute
%     \[
%         P_1^\infty \circ \pi_m^{-1}(A) 
%         = \int_{\probx} P(A) \, Q^1(dP) 
%         = P_0(A) 
%         = P_2^\infty \circ \pi_m^{-1}(A).
%     \]
% \end{remark}
% The implication of the above remark is the following: If $Q^1 \neq Q^2$, but $P_1^\infty\circ \pi_m^{-1}$ and  $P_2^\infty\circ \pi_m^{-1}$ are the same, then we observe samples drawn from the same probability measures. 
% In this case, no testing procedure can achieve a low Type II error, which is the probability of not rejecting the null hypothesis when it is false. To overcome this limitation, we need to guarantee that 
% \[
% Q^1 = Q^2 \quad \text{if and only if} \quad P^{1,(m)} = P^{2,(m)}.
% \]
% This indeed holds true for any $Q^1, Q^2$ for sufficiently large $m$ (see lemma \ref{lemma::finite_exchseq_equivalent_infinite_exchseq}).

% As noted above, hierarchical samples $\mathbf{X}_{n,m}\sim Q^1$ and $\mathbf{Y}_{n,m}\sim Q^2$ are used to decide whether $Q^1 = Q^2.$ We aim to develop a testing scheme inspired by the Kolmogorov–Smirnov (KS) test. Recall that, in the KS test, one decides to reject the null hypothesis if the distance between the empirical measures from the samples exceeds a given threshold. In section \ref{sec::distances}, we will equip the space $\probprobx$ with the distance function; for the moment, assume that $d$ defines a distance on that space. Then, we could naturally extend the idea of KS test to our setting if we had access to the samples $\widetilde{P}^1_1,\cdots, \widetilde{P}^1_n \overset{\iid}{\sim} Q^1$ and $\widetilde{P}^2_1,\cdots, \widetilde{P}^2_n \overset{\iid}{\sim} Q^2,$ by rejecting $H_0$ if 
% \[
% d\left( \frac{1}{n}\sum_{i = 1}^n \delta_{\widetilde{P}^1_i}, \frac{1}{n}\sum_{i = 1}^n \delta_{\widetilde{P}^2_i}\right)
% \]
% exceeds some threshold. But recall that we do not observe $\widetilde{P}^1_i$ and $\widetilde{P}^2_i$. Still, we can estimate them from the observations they generated, i.e. we estimate each of them using the corresponding sequences $X_{i,1},\ldots,X_{i,m}$ and $Y_{i,1},\ldots,Y_{i,m}$ by
% \[
% \widehat{P}^1_{i,m} = \frac{1}{m}\sum_{j = 1}^m\delta_{X_{i,j}}, \quad \widehat{P}^2_{i,m} = \frac{1}{m}\sum_{j = 1}^m\delta_{Y_{i,j}} \quad i = 1,\ldots, n.
% \]
% So, the idea of the test is to reject $H_0$ if 
% \[
% d\left( \frac{1}{n}\sum_{i = 1}^n \delta_{\widehat{P}^1_{i,m}}, \frac{1}{n}\sum_{i = 1}^n \delta_{\widehat{P}^2_{i,m}}\right)
% \]
% is higher than some threshold.

% Note that the law of each $\widehat{P}^k_{i,m}$ is different from $Q^k$ for $k = 1,2.$ Since such random probability measures are main objects to work with in this paper, we give a special notation to their laws. 

% In particular, given $Q \in \probprobx,$ we denote by $Q^{(m)}$ the law of empirical measure constructed from exchangeable sequence, i.e.,
% \[
% Q^{(m)} = \law \left( \frac{1}{m}\sum_{j = 1}^m \delta_{Z_j}\right),
% \]
% where
% $Z_i,\cdots,Z_m$ are exchangeable sequences from 
% \begin{align*}
% P &\ \sim Q \\
% Z_1, \ldots Z_m \mid P &\overset{\mathrm{i.i.d.}}{\sim} P.
% \end{align*}

% To sum up, given $Q^1, Q^2,$ since we estimate the random probability measures $\widehat{P}^k_{i,m}$ from the hierarchical samples $\mathbf{X}_{n,m}, \mathbf{Y}_{n,m}$ we say that we observe 
% \[
% \widehat{P}^1_{1,m},\cdots, \widehat{P}^1_{n,m} \overset{i.i.d}{\sim} Q^{1,(m)}, \quad \widehat{P}^2_{1,m},\cdots, \widehat{P}^2_{n,m} \overset{i.i.d}{\sim} Q^{2,(m)}.
% \]
% We want to decide whether $Q^1$ and $Q^2$ are the same by comparing the empirical measures concentrated on those two samples. Such empirical measures have a special name:
% \begin{definition}[Hierarchical estimator]
%     \label{def:hier_estimator}
%     Given $Q$, a the law of random probability measure, and hierarchical sample $\mathbf{X}_{n,m},$ the hierarchical estimator of $Q$ is defined as 
%     \[
%     Q_{n,m} = \frac{1}{n}\sum_{i = 1}^n \delta_{\widehat{P}_{i,m}}
%     \]
%     where $\widehat{P}_{i,m} = \frac{1}{m}\sum_{j = 1}^m \delta_{X_{i,j}}$ for all $i = 1,\ldots,n.$
% \end{definition}
% The testing scheme we develop will reject the null hypothesis whenever hierarchical estimators $Q^1_{n,m}$ and $Q^2_{n,m}$ are far away from each other. To clarify, what we mean by far away, we need a distance function on $\probprobx.$
